\documentclass[11pt]{article}
\usepackage{amsmath,amssymb,mathtools}
\usepackage{geometry}
\usepackage{bm}
\geometry{a4paper,margin=2.5cm}
\allowdisplaybreaks

\title{STEP1: Shock Propagation Along a Single Branch\\
Unifying the Second Version's Event Definitions with the First Version's Symmetric Formulas}
\author{}
\date{}

\begin{document}
\maketitle

\section{Purpose and Unified Conventions}

This document retains two complementary forms of representation:

\begin{enumerate}
    \item Use explicit event definitions to distinguish ``when a production node itself stops/recovers'' from
    ``when a specific downstream node actually observes that shock'';
    \item For a single shutdown window, preserve the highly symmetric
    stop/recovery recursions and cumulative formulas of the original version;
    \item For multiple windows and general DAGs, use the cumulative deficit function \(D(t)\)
    as a unified generalization of those formulas.
\end{enumerate}

The central principle is
\[
\boxed{
\text{Define physical events first, derive single-window formulas, then generalize to }D(t)\text{.}
}
\]

Consider a direct production-relation branch
\[
P_p\longrightarrow M_1\longrightarrow\cdots
\longrightarrow M_K\longrightarrow P_q.
\]

\section{The Two Types of Events Must Be Distinguished}

\subsection{The Production Node's Own Production-Start Event}

Let \(u_p(t)\) denote the actual intensity of production activity initiated by node \(P_p\) at time \(t\),
with initialized normal intensity
\[
u_p^0\equiv u_p(0).
\]

If \(P_p\) stops initiating new production during a single complete-shutdown window, define
\[
\boxed{
I_p^{act}=[\sigma_p,\rho_p),
}
\]
where
\[
\boxed{
\sigma_p
=
\text{time when \(P_p\) itself stops initiating new production},
}
\]
\[
\boxed{
\rho_p
=
\text{time when \(P_p\) itself resumes initiating production at normal intensity}.
}
\]

Define the duration of this internal shutdown window as
\[
\boxed{
\kappa_p=\rho_p-\sigma_p.
}
\]

Here \(\sigma_p,\rho_p\) describe only events on
\[
\boxed{\text{the production-start plane of \(P_p\)}}
\]
and do not imply that any particular downstream node has already experienced a supply interruption.

\subsection{Stop/Recovery Observed by a Specific Downstream Node}

For a specific direct downstream node \(P_q\), define
\[
\boxed{\begin{gathered}
t_{p\to q}^{\mathrm{stop}}=\\
\text{the first time \(P_q\) cannot obtain normally required inputs}\\
\text{because of the supply deficit from \(P_p\)}
\end{gathered}}
\]
and
\[
\boxed{\begin{gathered}
t_{p\to q}^{\mathrm{rec}}=\\
\text{the first time, after an effective interruption, that \(P_q\)}\\
\text{again receives normal supply from \(P_p\)'s resumed production}
\end{gathered}}
\]

Therefore, we must distinguish
\[
\boxed{
(\sigma_p,\rho_p)
\neq
(t_{p\to q}^{\mathrm{stop}},t_{p\to q}^{\mathrm{rec}}).
}
\]

The former is a
\[
\text{node-local event},
\]
The latter is a
\[
\text{downstream-visible event}.
\]

If \(P_p\) connects to multiple downstream nodes, then generally
\[
t_{p\to q_1}^{\mathrm{stop}}
\neq
t_{p\to q_2}^{\mathrm{stop}},
\qquad
t_{p\to q_1}^{\mathrm{rec}}
\neq
t_{p\to q_2}^{\mathrm{rec}},
\]
because different output branches may have different delays and inventory buffers.

\section{Two Basic Propagation Quantities}

For a direct branch \(P_p\to P_q\), let the transportation-edge and inventory-node sets be
\[
\mathcal E_{p\to q},
\qquad
\mathcal M_{p\to q}.
\]

\subsection{Total Time Shift}

Define
\[
\boxed{
\Theta_{p\to q}
=
\tau_p^{prod}
+
\sum_{e\in\mathcal E_{p\to q}}
\tau_e^{trans}.
}
\tag{1}
\]

Here, the upstream process \(P_p\)'s own production delay is explicitly included in
the downstream propagation from \(P_p\) to \(P_q\).

Therefore \(\Theta_{p\to q}\) represents only
\[
\boxed{\text{time shift}}
\]
and does not change the duration of a deficit.

\subsection{Total Inventory Coverage}

For each inventory node \(m\in\mathcal M_{p\to q}\),
denote its initialized normal outflow rate by
\[
f_m^{out,0}>0.
\]

Define inventory coverage
\[
c_m=\frac{q_m(0)}{f_m^{out,0}},
\]
The total inventory coverage time of the entire branch is then
\[
\boxed{
C_{p\to q}
=
\sum_{m\in\mathcal M_{p\to q}}
\frac{q_m(0)}{f_m^{out,0}}.
}
\tag{2}
\]

Physically:
\[
\boxed{
\text{Delay controls shock arrival; inventory controls remaining deficit duration.}
}
\]

\section{Main STEP1 Result: Symmetric Propagation Formulas for a Single Shutdown Window}

Let \(P_p\)'s own production-start gap be
\[
[\sigma_p,\rho_p),
\qquad
\kappa_p=\rho_p-\sigma_p.
\]

Define the candidate supply-interruption front
\[
\boxed{
s_{p\to q}
=
\sigma_p+\Theta_{p\to q}+C_{p\to q},
}
\tag{3}
\]
and the arrival front of resumed production output
\[
\boxed{
r_{p\to q}
=
\rho_p+\Theta_{p\to q}.
}
\tag{4}
\]

If
\[
s_{p\to q}<r_{p\to q},
\]
equivalently
\[
\boxed{
\kappa_p>C_{p\to q},
}
\tag{5}
\]
then this internal shutdown of \(P_p\) penetrates all inventories on the branch and affects \(P_q\).

In this case, the central STEP1 formulas are
\[
\boxed{
\begin{aligned}
t_{p\to q}^{\mathrm{stop}}
&=
\sigma_p
+\Theta_{p\to q}
+C_{p\to q},
\\[1mm]
t_{p\to q}^{\mathrm{rec}}
&=
\rho_p
+\Theta_{p\to q}.
\end{aligned}
}
\tag{6}
\]

This preserves the most direct symmetric structure of the original version:
\[
\boxed{
\begin{aligned}
\text{stop}
&=
\text{upstream stop}
+\text{delay}
+\text{inventory coverage},
\\
\text{recovery}
&=
\text{upstream recovery}
+\text{delay}.
\end{aligned}
}
\tag{7}
\]

Subtracting the two formulas yields
\[
\boxed{
t_{p\to q}^{\mathrm{rec}}
-
t_{p\to q}^{\mathrm{stop}}
=
(\rho_p-\sigma_p)-C_{p\to q}.
}
\tag{8}
\]

Thus the effective duration of the supply interruption is
\[
\boxed{
\Delta_{p\to q}
=
\left[
(\rho_p-\sigma_p)-C_{p\to q}
\right]_+.
}
\tag{9}
\]

If
\[
s_{p\to q}\ge r_{p\to q},
\]
then
\[
\boxed{
\Delta_{p\to q}=0.
}
\tag{10}
\]

In this case, although \(P_p\) itself did indeed shut down,
the resumed-output front catches up to the deficit front before the inventories are exhausted,
so \(P_q\) never experiences a genuine effective supply interruption.
Here \(r_{p\to q}\) can be called the ``resumed-output arrival front,''
but it should not be described as \(P_q\) ``shutting down and then recovering.''

\section{Standard \(P-M-M-P\) Structure}

Consider
\[
P_p\rightarrow M_1\rightarrow M_2\rightarrow P_q.
\]

Define
\[
\boxed{
\begin{aligned}
\Theta_{p\to q}
={}&
\tau_p^{prod}
+\tau_{P_p\to M_1}^{trans}
+\tau_{M_1\to M_2}^{trans}
+\tau_{M_2\to P_q}^{trans},
\\[1mm]
C_{p\to q}
={}&
\frac{q_{M_1}(0)}
     {f_{M_1\to M_2}^{0}}
+
\frac{q_{M_2}(0)}
     {f_{M_2\to P_q}^{0}}.
\end{aligned}
}
\tag{11}
\]

If a single shutdown window penetrates all inventories, then
\[
\boxed{
\begin{aligned}
t_{p\to q}^{\mathrm{stop}}
={}&
\sigma_p
+\tau_p^{prod}
+\tau_{P_p\to M_1}^{trans}
+\frac{q_{M_1}(0)}{f_{M_1\to M_2}^{0}}
\\
&+
\tau_{M_1\to M_2}^{trans}
+\frac{q_{M_2}(0)}{f_{M_2\to P_q}^{0}}
+\tau_{M_2\to P_q}^{trans},
\\[2mm]
t_{p\to q}^{\mathrm{rec}}
={}&
\rho_p
+\tau_p^{prod}
+\tau_{P_p\to M_1}^{trans}
+\tau_{M_1\to M_2}^{trans}
+\tau_{M_2\to P_q}^{trans}.
\end{aligned}
}
\tag{12}
\]

Thus, production and transportation delays
appear perfectly symmetrically in the stop and recovery fronts;
inventory coverage enters only the stop front.

\section{Single-Source Linear Chain: Recovering the Original Recursive and Cumulative Forms}

Consider
\[
P_0
\rightarrow\mathcal B_1
\rightarrow P_1
\rightarrow\mathcal B_2
\rightarrow P_2
\rightarrow\cdots
\rightarrow\mathcal B_i
\rightarrow P_i,
\]
where \(\mathcal B_j\) denotes
a material branch from \(P_{j-1}\) to \(P_j\) containing no other production nodes.

Define
\[
\Theta_j
\equiv
\Theta_{P_{j-1}\to P_j},
\qquad
C_j
\equiv
C_{P_{j-1}\to P_j}.
\]

Suppose the source \(P_0\) experiences a single production-start gap
\[
[\sigma_0,\rho_0).
\]

In a single-input linear chain, as long as the shock has not been fully absorbed by inventory,
the downstream-visible shortage from \(P_{j-1}\) to \(P_j\)
directly determines \(P_j\)'s own production-start gap.
Hence the recursion is
\[
\boxed{
\begin{aligned}
\sigma_j
&=
\sigma_{j-1}
+\Theta_j
+C_j,
\\[1mm]
\rho_j
&=
\rho_{j-1}
+\Theta_j.
\end{aligned}
}
\tag{13}
\]

Summing segment by segment gives
\[
\boxed{
\begin{aligned}
\sigma_i
&=
\sigma_0
+\sum_{j=1}^{i}\Theta_j
+\sum_{j=1}^{i}C_j,
\\[1mm]
\rho_i
&=
\rho_0
+\sum_{j=1}^{i}\Theta_j.
\end{aligned}
}
\tag{14}
\]

Therefore
\[
\boxed{
\rho_i-\sigma_i
=
(\rho_0-\sigma_0)
-
\sum_{j=1}^{i}C_j.
}
\tag{15}
\]

If the earlier version assumes
\[
\sigma_0=0,
\qquad
\rho_0=k,
\]
we immediately obtain
\[
\boxed{
\begin{aligned}
\sigma_i
&=
\sum_{j=1}^{i}(\Theta_j+C_j),
\\[1mm]
\rho_i
&=
k+\sum_{j=1}^{i}\Theta_j,
\end{aligned}
}
\tag{16}
\]
and
\[
\boxed{
\rho_i-\sigma_i
=
k-\sum_{j=1}^{i}C_j.
}
\tag{17}
\]

If
\[
k\le \sum_{j=1}^{i}C_j,
\]
then the shock has already been fully absorbed by accumulated inventory coverage before reaching \(P_i\),
so \(P_i\) develops no effective production-start gap.

\section{From the Original Cumulative Formulas to the Cumulative Deficit Function \(D(t)\)}

The original symmetric formulas are suitable for describing
\[
\boxed{\text{a single stop/recovery window}.}
\]

To handle multiple disjoint windows, branching after merging, and general DAGs,
define the cumulative internal production-start deficit of \(P_p\) as
\[
\boxed{
D_p^{act}(t)
=
\text{Equivalent normal-production time not initiated by \(P_p\) up to time \(t\)}.
}
\tag{18}
\]

Under the complete-shutdown / complete-recovery baseline,
\[
\dot D_p^{act}(t)\in\{0,1\}
\]
holds almost everywhere, and
\[
\boxed{
u_p(t)
=
u_p^0\left[1-\dot D_p^{act}(t)\right].
}
\tag{19}
\]

For a single internal shutdown window
\[
[\sigma_p,\rho_p),
\]
we have
\[
\boxed{
D_p^{act}(t)
=
\min\left\{
(t-\sigma_p)_+,
\rho_p-\sigma_p
\right\}.
}
\tag{20}
\]

Adopt the convention that when \(s<0\),
\[
D_p^{act}(s)=0.
\]

Then the STEP1 propagation operator for arbitrary windows is defined as
\[
\boxed{
D_{p\to q}(t)
=
\mathcal T_{p\to q}[D_p^{act}](t)
=
\left[
D_p^{act}(t-\Theta_{p\to q})
-
C_{p\to q}
\right]_+.
}
\tag{21}
\]

Here
\[
(x)_+\equiv\max\{x,0\}.
\]

Its basic components are
\[
\boxed{
\text{Delay}:D(t)\mapsto D(t-\tau),
}
\tag{22}
\]
and
\[
\boxed{
\text{Inventory}:D(t)\mapsto
\left[D(t)-\frac{q(0)}{f^0}\right]_+.
}
\tag{23}
\]

Since
\[
[[x-c_1]_+-c_2]_+
=
[x-c_1-c_2]_+,
\]
the coverage times of multiple sequential inventories can be added directly.

\section{How the \(D(t)\) Operator Reduces to the Original Symmetric Formulas}

Substitute the single-window equation \((20)\) into \((21)\):
\[
D_{p\to q}(t)
=
\left[
\min\left\{
(t-\Theta_{p\to q}-\sigma_p)_+,
\rho_p-\sigma_p
\right\}
-
C_{p\to q}
\right]_+.
\tag{24}
\]

If
\[
\rho_p-\sigma_p>C_{p\to q},
\]
then the time when \(D_{p\to q}(t)\) begins to increase is precisely
\[
\boxed{
t_{p\to q}^{\mathrm{stop}}
=
\sigma_p+\Theta_{p\to q}+C_{p\to q},
}
\]
and the time when it stops increasing is precisely
\[
\boxed{
t_{p\to q}^{\mathrm{rec}}
=
\rho_p+\Theta_{p\to q}.
}
\]

Therefore
\[
\boxed{\begin{gathered}
\text{\(D(t)\) is not a separate model;}\\
\text{it generalizes the original stop/recovery formulas to multiple windows.}
\end{gathered}}
\tag{25}
\]

\section{Operator Composition and the Original ``Segment-by-Segment Summation''}

Define
\[
\mathcal T_{\Theta,C}[D](t)
=
[D(t-\Theta)-C]_+.
\]

Across two branch segments:
\[
D_1(t)
=
\mathcal T_{\Theta_1,C_1}[D_0](t),
\]
\[
D_2(t)
=
\mathcal T_{\Theta_2,C_2}[D_1](t).
\]

Using
\[
[[x-C_1]_+-C_2]_+
=
[x-C_1-C_2]_+,
\]
we obtain
\[
\boxed{
D_2(t)
=
\left[
D_0(t-\Theta_1-\Theta_2)
-(C_1+C_2)
\right]_+.
}
\tag{26}
\]

Therefore, repeated composition along a single-source linear chain gives
\[
\boxed{
D_i^{act}(t)
=
\left[
D_0^{act}
\left(
t-\sum_{j=1}^{i}\Theta_j
\right)
-
\sum_{j=1}^{i}C_j
\right]_+.
}
\tag{27}
\]

Thus, on a chain without merges, the new operator
agrees exactly with the original rule of summing delays and inventory coverage segment by segment.

\section{Correspondence with Daily Production Status for Plant \(i\) on Day \(k\)}

\(D_i^{act}(t)\) is the continuous-time cumulative representation of the original daily production states.

Under the binary baseline,
\[
\boxed{
\frac{u_i(t)}{u_i^0}
=
1-\dot D_i^{act}(t).
}
\tag{28}
\]

Therefore
\[
\dot D_i^{act}(t)=0
\quad\Longleftrightarrow\quad
P_i\text{ produces normally},
\]
\[
\dot D_i^{act}(t)=1
\quad\Longleftrightarrow\quad
P_i\text{ is completely shut down}.
\]

If day \(k\) is defined as
\[
[k,k+1),
\]
and no additional intraday switching takes place, we can write
\[
\boxed{
x_i[k]
=
1-
\left(
D_i^{act}(k+1)-D_i^{act}(k)
\right),
}
\tag{29}
\]
where
\[
x_i[k]\in\{0,1\}.
\]

Then
\[
\boxed{
x_i[k]=1
\Longleftrightarrow
\text{Plant \(i\) operates on day \(k\)},
}
\]
\[
\boxed{
x_i[k]=0
\Longleftrightarrow
\text{Plant \(i\) is shut down on day \(k\)}.
}
\]

If stop/recovery switching is allowed within a day, then
\[
D_i^{act}(k+1)-D_i^{act}(k)\in[0,1]
\]
naturally represents the proportion of normal production time lost on that day.

\section{Defining Stop/Recovery for Multiple Windows}

If \(P_p\) experiences multiple disjoint production-start gaps,
the same downstream inventory must not be reinitialized to \(q(0)\) for each window.

Thus one cannot apply
\[
C_{p\to q}
\]
separately to each window and simply take the union of the resulting intervals.

The correct procedure is first to calculate the entire
\[
D_{p\to q}(t)
=
[D_p^{act}(t-\Theta_{p\to q})-C_{p\to q}]_+.
\]

Under the complete-shutdown / complete-recovery baseline,
all maximal connected intervals satisfying
\[
\dot D_{p\to q}(t)=1
\]
are defined as
\[
\boxed{
[t_{p\to q}^{\mathrm{stop},\ell},
t_{p\to q}^{\mathrm{rec},\ell}),
\qquad
\ell=1,\ldots,L_{p\to q}.
}
\tag{30}
\]

Thus, in the general case,
\[
\boxed{
D_{p\to q}(t)
\text{ is the primary interface,}
\qquad
(t^{\mathrm{stop}},t^{\mathrm{rec}})
\text{ is an event label read from \(D\).}
}
\]

\section{Normal Flow Rates and Fixed Proportions}

If the output coefficient of \(P_p\) for some product \(Y\) is
\[
b_{pY}>0,
\]
and a fixed fraction
\[
\rho_{p\to q}>0
\]
is allocated to the branch leading to \(P_q\), with nonbinding output-edge capacity,
then the initialized normal flow rate on that branch is
\[
\boxed{
f_{p\to q}^0
=
\rho_{p\to q}b_{pY}u_p^0.
}
\tag{31}
\]

Fixed output coefficients and fixed allocation ratios only rescale material quantities.
After normalizing deficits by their corresponding normal rates into ``equivalent normal-production time,''
they do not alter the shape of \(D(t)\);
they enter inventory coverage through the normal flow rates
\[
q(0)/f^0.
\]

\section{Applicability Limits of the Current Rules}

The foregoing closed-form expressions and operators rely on the following baseline:
\begin{itemize}
    \item The network is initialized at steady state for \(t<0\);
    \item Production and transportation delays are fixed;
    \item Production recipes and output allocation ratios are fixed;
    \item Output-edge capacity is nonbinding;
    \item No active restocking, overproduction, or rerouting;
    \item When normal flow resumes into a depleted inventory, it may pass through immediately without first replenishing the stock to \(q_m(0)\);
    \item No losses, scrapping, or shelf-life constraints;
    \item The main focus is binary production intensity: complete shutdown or complete recovery.
\end{itemize}

If partial recovery is allowed later,
the cumulative deficit function still applies, but generally
\[
0\le \dot D(t)\le1,
\]
and
\[
u_p(t)=u_p^0[1-\dot D_p^{act}(t)]
\]
can describe partial production.

If an equipment failure directly freezes WIP already inside the process,
production delay is no longer a simple pure delay,
and an explicit WIP state is required; this case is not considered here.

\section{Final STEP1 Interface}

For a single window, use the intuitive symmetric formulas first
\[
\boxed{
\begin{aligned}
t_{p\to q}^{\mathrm{stop}}
&=
\sigma_p+\Theta_{p\to q}+C_{p\to q},
\\
t_{p\to q}^{\mathrm{rec}}
&=
\rho_p+\Theta_{p\to q},
\end{aligned}
}
\]
and check
\[
\rho_p-\sigma_p>C_{p\to q}.
\]

For multiple windows or general DAGs, use the unified operator
\[
\boxed{
D_{p\to q}(t)
=
\left[
D_p^{act}(t-\Theta_{p\to q})
-C_{p\to q}
\right]_+.
}
\]

STEP1 can therefore be summarized as
\[
\boxed{
\underbrace{D_p^{act}}_{\text{upstream node's own activity}}
\quad
\xrightarrow{
\substack{
\text{production/transport delay}\\
\text{inventory absorption}
}}
\quad
\underbrace{D_{p\to q}}_{\text{shortage visible to a specific downstream node}}.
}
\]

This \(D_{p\to q}(t)\) is exactly the input received directly by STEP2 during its required-input merge.

\end{document}
